\section{Conditional Trajectory Abstraction}
\label{sec:method}

\method has three learned parts: a target encoder that turns the actual future into a code (\cref{sec:source}), two heads that read the code (\cref{sec:reader}), and a world model that predicts the code from an action (\cref{sec:wm}).

\subsection{Inputs}
\label{sec:context}
As in DINO-WM~\cite{zhou2025dinowm}, images are embedded by a frozen DINOv2 ViT-S/14~\cite{oquab2024dinov2}; we project the patch features to 128 dimensions with PCA.
The context $C$ contains the $16{\times}16$ patch tokens of the current image, the previous image pooled to $8{\times}8$, and the proprioception at both times.
The future segment $\tau$ contains the full token grid of the final image, pooled tokens of intermediate images with temporal embeddings, and the final proprioception.
Features are computed once per image and shared by all candidates and goals.

\subsection{Target encoder}
\label{sec:source}
The target encoder $q_\phi(S\mid C,\tau)$ is used only during training.
$M$ learned queries cross-attend to the future and context tokens, and each output is quantized by FSQ~\cite{mentzer2024fsq} with levels $(8,8,4)$, giving 256 values per token.
Because the queries see the context, the encoder learns which parts of the future the history already implies and spends the code on the rest.
The $M$ tokens describe the segment as a whole, not one frame each.

\subsection{Reader and feature decoder}
\label{sec:reader}
The reader $D_\psi(C,S,g)$ is a Transformer over the context tokens, the code tokens, and the patch tokens of a goal image, and outputs one score.
It is trained on candidates that branch from the same state, with a pairwise logistic loss over pairs whose labels differ by more than a margin $\epsilon$:
\begin{equation}
\begin{aligned}
  \mathcal{L}_{\mathrm{rank}}&=\E_{(i,j):\,\ell_i>\ell_j+\epsilon}\big[\log\big(1+e^{-(s_i-s_j)}\big)\big],\\
  s_i&=D_\psi(C,S_i,g).
\end{aligned}
\label{eq:rank}
\end{equation}
Ranking within a decision matches how the score is used, so it need not be calibrated across states.

A ranking loss alone would let the code shrink to whatever orders the training goals.
The feature decoder $G_\psi(C,S)$ prevents this by reconstructing the future tokens from the context and the code.
It predicts each token's change from the current image, and its loss $\mathcal{L}_{\mathrm{feat}}$ is normalized by that of predicting no change, so an uninformative code scores one.
In JEPA terms, the decoder keeps the target representation from collapsing~\cite{lecun2022path}; no image is ever generated.

\subsection{World model}
\label{sec:wm}
The world model $p_\theta(S\mid C,\bA)$ is the JEPA predictor.
A Transformer encoder reads the context tokens and one token per action, and $M$ learned queries decode all code tokens in parallel.
Each FSQ coordinate of each token is predicted as a categorical distribution over its levels, and the reader receives the expected codeword $\hat S$.
A single forward pass therefore scores the whole set of proposals.
The world model never sees future images, the goal, or labels at test time.

It is trained against the stop-gradient code $S=\sg(q_\phi(C,\tau))$:
\begin{equation}
  \mathcal{L}_{\mathrm{WM}}=\mathcal{L}_{\mathrm{NLL}}+\mathcal{L}_{\mathrm{cons}}+\mathcal{L}^{\hat S}_{\mathrm{rank}}.
  \label{eq:wm}
\end{equation}
$\mathcal{L}_{\mathrm{NLL}}$ is the cross-entropy of each coordinate of $S$.
The other two terms pass through the frozen reader: $\mathcal{L}_{\mathrm{cons}}$ matches the within-decision score differences of $\hat S$ to those of $S$, and $\mathcal{L}^{\hat S}_{\mathrm{rank}}$ applies~\eqref{eq:rank} to the predicted codes, weighting pairs by their label difference.
A context-only model $r_\omega(S\mid C)$ is trained with $\mathcal{L}_{\mathrm{NLL}}$ alongside; its excess cross-entropy over the world model estimates $I(S;\bA\mid C)$, the part of the code the action explains.

\subsection{Training}
\label{sec:codesign}
\paragraph{Branched data.}
At every decision of a policy rollout, we execute each proposal for $\He$ steps from a clone of the simulator state, obtaining several futures and labels from one history.
To make siblings differ enough to rank, we also branch a copy of each proposal shifted by a constant offset.
In half of the rollouts, half of the decisions follow the best branch rather than the default one, so that the data cover the states a good selector visits.
All branches of an episode stay in one split.

\paragraph{Stages.}
Stage 1 trains the target encoder, reader, and decoder on actual futures with $\mathcal{L}_{\mathrm{code}}=\mathcal{L}_{\mathrm{rank}}+\alpha\,\mathcal{L}_{\mathrm{feat}}$, using straight-through gradients through the quantizer.
Stage 2 freezes them and trains the world model with~\eqref{eq:wm}; each network keeps the checkpoint with the best development objective.
An optional stage co-designs the code for predictability, alternating code and world-model updates with a penalty $\lambda\lVert S-\E_{\bar p_\theta}[S]\rVert^2$ under a frozen copy of the world model; we treat it as an ablation.
Throughout, we monitor code perplexity, distinct codes per decision, the decoder's $R^2$, and $I(S;\bA\mid C)$ to detect collapse.

\subsection{Planning}
\label{sec:planning}
\Cref{alg:cta} summarizes deployment.
Per decision, \method encodes the images once, runs the world model once on all proposals, and runs the reader on every candidate--goal pair; a new goal costs only reader evaluations.

\begin{algorithm}[t]
\caption{Selecting policy proposals with \method.}
\label{alg:cta}
\small
\begin{algorithmic}[1]
\Require frozen policy $\pi$ and encoder $\varphi$; world model $p_\theta$; reader $D_\psi$; goal images $\cG$
\While{episode not finished}
  \State $C\gets\varphi(h_t)$ \Comment{once per decision}
  \State $\{\bA^{(k)}\}_{k=1}^{K}\gets$ proposals of $\pi$
  \State $\hat S^{(k)}\gets$ expected code of $p_\theta(\cdot\mid C,\bA^{(k)})$ \Comment{one batched pass}
  \State $s_k\gets\frac{1}{|\cG|}\sum_{g\in\cG}D_\psi\big(C,\hat S^{(k)},g\big)$
  \State $k^\star\gets$ smallest index attaining $\max_k s_k$ \Comment{ties keep $\bA^{(1)}$}
  \State execute $\bA^{(k^\star)}_{1:\He}$; $t\gets t+\He$
\EndWhile
\end{algorithmic}
\end{algorithm}

\subsection{Implementation details}
\label{sec:impl}
All modules have width 256 and 8 heads: the target encoder has three decoder layers, the reader four encoder layers, the feature decoder two decoder layers, and the world model four encoder and four decoder layers.
We use $M{=}16$ (128 bits), $\alpha{=}1$, and $\epsilon{=}10^{-3}$.
Stage 2 runs 8{,}000 AdamW updates (learning rate $10^{-4}$, weight decay $0.05$, bfloat16) with 32 decisions per update.
\TODO{Final stage-1 schedule.}
On PushT, $96{\times}96$ images are resized to $224{\times}224$ for DINOv2, $\He{=}8$, and the intermediate images follow steps 2, 4, and 6.
Each action is a target position of the agent, embedded in absolute and relative coordinates.
